\documentclass[11pt]{article}
\usepackage{amsmath,amssymb,amsthm}
\newtheorem{proposition}{Proposition}
\title{An exact interior-window lift: integrate the boundary-adjacent bath first}\author{}\date{}
\begin{document}\maketitle
\paragraph{Original law, not the previous whole-path projection inequality.}
Let $X$ be the entire initial spatial slice and $Y$ the midpoint slice of the
original periodic Wilson slab of temporal length $2\ell$. Its true marginal
is $p(dX)$. Retain a plaquette-disjoint set $I$ of spatial-link trajectories,
$N_I=|I|$, as in the one-color argument. The original purity identity is
$\mathcal P_{\rm full}=E_{X,X'\ iid\ p}\operatorname{Aff}(p_X^Y,p_{X'}^Y)^2$.

Choose a window $[\ell-w,\ell+w]$ with $w<\ell$, disjoint from the initial
slice. Define $\Omega$ to contain all complementary spatial and temporal links
needed in plaquettes touching the selected trajectories inside this window,
and all complementary midpoint spatial links. For simplicity one may retain
all complementary spatial slices and temporal links in that window, including
its two bounding collars. No boundary-adjacent path outside the window is
retained. Let
\[
 \nu_X(d\omega)=\mathbb P(d\Omega=\omega\mid X)
\]
be its EXACT marginal under the original Wilson law. All outside variables,
including boundary-adjacent bath paths and selected variables, have already
been integrated in this definition; no conditional chain normalizer is dropped.
Let $\alpha=(u_i(\ell-w),u_i(\ell+w))_{i\in I}$ be the selected endpoint vector.

\paragraph{Exact finite-range disintegration.}
Conditional on $(X,\Omega=\omega,\alpha)$, the selected paths strictly inside
the window have the independent original conditional-chain laws, with endpoints
$\alpha$ and the literal bath $\omega$. To see this, assign every original
plaquette containing a selected interior link to its selected chain. All its
other links lie in the retained collar record or in that same chain. Distinct
selected chains share no plaquette. Every remaining action term is independent
of all selected interior links; it can affect only the posterior of $\omega$
and $\alpha$. Integrating that remainder therefore gives the exact identity
\[
 p_X^{U\mid\omega}=\int\sigma_X(d\alpha\mid\omega)
                              \bigotimes_{i\in I}q_{i,\omega}^{\alpha_i}, \tag{1}
\]
where $U=(u_i(\ell))_{i\in I}$ and $\sigma_X$ is the ACTUAL endpoint posterior.
It can be arbitrarily correlated across $i$ and depend arbitrarily on $X$.
The normalized chain laws $q$ in(1) are independent of $X$ once their endpoints
and the bath window are fixed. Their normalization factors are part of the
posterior $\sigma_X$ and marginal $\nu_X$, not canceled out of the full measure.
This is the new disintegration needed to shrink the lift; it does not follow
by simply projecting the previous whole-path bath affinity.

\paragraph{Good runs remove endpoint dependence before taking mixtures.}
Use the actual bath criterion and sparse-record parameter of the original sparse bath records calculation,
\[
 p_B=6K_t^{1/12}\exp[-tB(1-\cos\epsilon_0)/108].
\]
For $L\ge2$, $h=L+4$, and $\delta=3e^{64B}(2/5)^L\le1/2$, let $\mathcal G$
be the window-measurable event that each selected trajectory has an all-good
candidate block of $h$ labels on EACH side of the midpoint. Candidate blocks
are chosen inside the collars so that their inner $L+2$ steps have all onsite
and temporal labels present. The good-run argument of the good run arbitrary complement calculation gives, uniformly
in arbitrary pairs of endpoint values,
\[
 \operatorname{Aff}(q_{i,\omega}^{\alpha_i},
                    q_{i,\omega}^{\beta_i})^2\ge1-16\delta^2,
 \qquad\omega\in\mathcal G.                                            \tag{2}
\]
Only original positive kernels lie outside the good runs, including the
rest of the window. Thus no endpoint or exterior bound is assumed here.

Affinities multiply for the product laws in(1). They are jointly concave under
ANY coupling of mixing distributions: pointwise Cauchy--Schwarz gives
$\sqrt{(\int f\,d\pi)(\int g\,d\pi)}\ge\int\sqrt{fg}\,d\pi$.
Use $\pi=\sigma_X\otimes\sigma_{X'}$, without assuming the components of
$\alpha$ independent. Equations(1)--(2) yield
\[
 \operatorname{Aff}(p_X^{U\mid\omega},p_{X'}^{U\mid\omega})
       \ge a_I:=(1-16\delta^2)^{N_I/2},\qquad\omega\in\mathcal G.         \tag{3}
\]
Assume $16\delta^2\le1$, as holds in the application below.

\paragraph{A new, genuinely thinner bath factor.}
The pair $(\Omega,U)$ projects to the full midpoint $Y$, since the complementary
midpoint links were retained in $\Omega$. Integrating conditional affinities
and then applying this projection gives
\[
 \operatorname{Aff}(p_X^Y,p_{X'}^Y)
   \ge a_I\int_{\mathcal G}\sqrt{d\nu_Xd\nu_{X'}}.                       \tag{4}
\]
Define the actual window factor
\[
 \mathcal P_{\rm win}=E_{X,X'\ iid\ p}\operatorname{Aff}(\nu_X,\nu_{X'})^2.
\]
Averaged excluded overlap is at most the ORIGINAL probability of the bad window,
by Fubini and Cauchy--Schwarz. Thus the same elementary square estimate as in
the good run arbitrary complement calculation now proves the NEW inequality
\[
 \mathcal P_{\rm full}\ge(1-16\delta^2)^{N_I}
                  [\mathcal P_{\rm win}-2\mathbb P(\mathcal G^c)].       \tag{5}
\]
The whole-path bath factor has been replaced by the actual window factor using
(1)--(3), not by reversing the direction of a projection inequality.

\paragraph{Original probability and a vanishing physical window thickness.}
Put $J$ disjoint candidate blocks on each side, so $w=Jh+O(1)$ suffices.
If no block for one side is entirely good, choose a bad label from each block.
The original subset bound, followed by a union bound over the two sides and
$N_I$ trajectories, gives
\[
 \mathbb P(\mathcal G^c)\le2N_I(hp_B)^J.                               \tag{6}
\]
This is an unconditional bound, precisely the one required in(5). No uniform
estimate conditional on a fixed rough $X$ has been introduced.

For a concrete choice let
\[
 L=\max\{2,\lceil[64B+10\log(2+\ell)]/\log(5/2)\rceil\},\qquad
 J=\max\{1,\lceil\log[2N_I(2+\ell)^{10}]/\log[1/(hp_B)]\rceil\},       \tag{7}
\]
when $hp_B<1$. Then $\delta\le3(2+\ell)^{-10}$ and
$\mathbb P(\mathcal G^c)\le(2+\ell)^{-10}$.
On the original regime $B\sim b\log\ell$, fixed $b>0$, bounded $K_t$, and
polynomial $N_I$ in $\ell$, the denominator defining $J$ is of order
$\log\ell$. Hence $J=O(1)$, $w=O(\log\ell)$, and $w<\ell$ eventually.
At fixed physical half-slab time, $\ell\asymp a^{-1}$, this window has physical
thickness $O(a\log(1/a))\to0$. Both errors in(5) vanish for the usual
$N_I=O(a^{-3})$ spatial count.

\paragraph{What has and has not been reduced.}
The retained record no longer contains paths near the initial boundary.
It contains a thin interior bath window with original posterior density,
including all induced interactions, all normalizers and global sectors.
Its affinity is generally smaller than the affinity of its projection to the
remaining midpoint slice. Consequently $\mathcal P_{\rm win}$ is STILL NOT
an output midpoint purity. A new continuity/coarse-response estimate would
be required to discard the remaining window information, or a spatial
iteration must retain this actual window law. Neither follows merely because
its physical thickness tends to zero. Equation(5) is a constructive reduction
of the former overly informative lift, not a lower bound for
$\mathcal P_{\rm win}$ or a full S94/mass-gap proof.

\paragraph{The second color induced interaction calculation and the next step.}
The exact output density $e^{-S_{\rm rest}}\prod_i\mathcal Z_i/Z$, its
conditional covariance Hessian, the cancelled-staple example with induced
mixed derivative $-B^2/4$, and the original-law $L^p$ bound $C_pB$ are
verified directly. The last bound uses $|\partial s|^2\le2s$, same-Wilson moments,
conditional Cauchy--Schwarz and outer H\"older. The two-arc averaged temporal
decay estimate also holds, with no uniform rough-boundary assertion. None
of these alone identifies the window law here as a Wilson or independent-chain
law for a second spatial color.
\end{document}
